% ---------------------------------
\section{Introduction}

\subsection{Background: Neeli Prasad's Seminar on Human-AI Collaboration}

Neeli Prasad's IEEE Distinguished Lecture walked through three phases of AI: perception AI (medical imaging, autonomous vision systems), generative AI (large language models, code generation), and physical AI (humanoid robots). She said ethics, privacy, and governance must be built in from the start \cite{prasad2026seminar}. Privacy especially. It's not something you add later.

She gave examples of what happens when privacy fails. A Hong Kong bank lost millions in a deepfake fraud where AI impersonated a CEO's voice and video. Patient records have been leaked and used for blackmail. Images and audio of minors have been faked to harass them, linked to documented suicides. Without privacy protection, AI causes real damage.

\subsection{Why Privacy in AI Systems}

The topic feels urgent right now. The EU updated the AI Act in 2024 \cite{eu_ai_act_2024}. GDPR was 2016 \cite{gdpr2016}. That ten-year lag suggests AI privacy law is moving faster now. Anyone entering engineering should know this. It's not optional.

Prasad also pushed for bringing different disciplines together. Privacy isn't just a technical problem. You need engineers, lawyers, ethicists, psychologists in the room at once, thinking about consent, data ownership, surveillance, labor displacement, and human dignity.

\subsection{Project Objective}

This project takes a practical approach. Rather than write an essay, it runs privacy research through four AI tools in sequence. Perplexity gathers research. Claude synthesizes it. QuillBot refines the writing. SlidesAI.io generates slides. The whole process is documented, including every prompt and output \cite{prasad2026seminar}. This follows Prasad's insistence that students show their prompts and reasoning.

\subsection{The Four-Tool Pipeline}

The pipeline works in four stages:

\begin{enumerate}
    \item \textbf{Perplexity} (Research): Gather cited information on AI privacy risks, regulatory frameworks, and real-world cases.
    \item \textbf{Claude API} (Synthesis): Analyze the seminar content and research findings, then generate a comprehensive article on privacy mechanisms, compliance challenges, and technical safeguards.
    \item \textbf{QuillBot} (Refinement): Rewrite content in different styles to see how tone and wording affect clarity.
    \item \textbf{SlidesAI.io} (Presentation): Turn the curated content into a slide deck.
\end{enumerate}

\subsection{Document Structure}

Section 2 covers the detailed prompts, settings, and execution for each tool. Section 3 shows the actual outputs and compares their quality and efficiency. Section 4 concludes with what was learned and connects back to Prasad's themes: transparency, governance, and responsible human-AI collaboration.
